\section{Data and Evaluation Setup}
\label{sec:data}

\paragraph{Concepts and prompts.}
We adopt the fear-concept dataset used by prior activation-steering
work~\citep{liu2024iti} as adapted in our codebase. From the upstream concept
list we sample five fear concepts under seed $7$:
\textit{Bugs}, \textit{Enclosed spaces}, \textit{Strangers}, \textit{Shadows},
\textit{Men}. Each concept has an associated paired statement pool, where every
positive statement is the same as the corresponding negative statement except
for the prepended one-line concept prefix (\emph{``Personify someone who is
terrified of \textit{<concept>}.''}). These paired statements drive the
routing-trace collection (Section~\ref{sec:methods}).

\paragraph{Evaluation prompts.}
For evaluation we sample five questions per concept under the same seed, drawn
from the upstream phobia-evaluation prompt family. The five questions are
generic enough to elicit fear-themed answers without naming the concept:

\begin{itemize}
  \item v1: \textit{What is the scariest thing in the world? Answer in 20
  words or less.}
  \item v2: \textit{Is there anything you secretly hope never happens to you?
  Answer in 20 words or less.}
  \item v3: \textit{What would a movie villain based on your worst fear be
  like? Answer in 20 words or less.}
  \item v4: \textit{If a fear had to visit you in your sleep, what form would
  it take? Answer in 20 words or less.}
  \item v5: \textit{What would a fear-themed amusement park ride, made just
  for you, be like? Answer in 20 words or less.}
\end{itemize}

This gives $5 \times 5 = 25$ comparison cases per condition. The
$20$-word cap is part of the prompt, which lets us treat
``response exceeds 20 words'' as a coarse coherence signal.

\paragraph{The two prompt contracts.}
The same paired statement pool feeds two distinct evaluation contracts:

\begin{description}
  \item[Prefix-conditioned (model-suitability).] At test time the model sees
  \texttt{<concept prefix> + <evaluation question>}. This is the contract for
  Claim~1 and is shared by the \llama, \olmoe, and \mixtral~baselines.
  \item[Question-only (steering test).] At test time the model sees only
  \texttt{<evaluation question>}. The concept prefix is intentionally
  omitted. This is the contract for Claims~2 and~3, and the
  \steermoe-Mixtral column adds a router-logit bias derived from the
  prefix-conditioned vs.\ unprefixed routing traces. The Llama-Mixtral
  reference column is unsteered.
\end{description}

\paragraph{What we record per case.}
For every (concept, evaluation question, condition) triple we record the
generated response. From each response we derive three lightweight diagnostic
statistics --- word count, concept-token hit (does any token of the concept
name appear in the response, after lowercasing), and a $20$-word coherence
flag --- and report aggregates over the $25$ cases per (condition, run).
Section~\ref{sec:methods} describes the routing-trace and steering-plan
artifacts that drive the SteerMoE column.

\paragraph{Reproducibility.}
All sampled concepts, sampled questions, paired statement bodies, and
generation metadata are committed under
\texttt{experiments/<run\_dir>/manual\_review\_plan.json} and
\texttt{generation\_metadata.json}. The aggregation script
\texttt{paper/scripts/aggregate\_results.py} re-derives the paper's tables
from those committed artifacts; it does not load any model weights.
